\hypertarget{ux5b89ux88c5ux7b2cux4e09ux65b9ux6a21ux5757}{%
\subsection{安装第三方模块}\label{ux5b89ux88c5ux7b2cux4e09ux65b9ux6a21ux5757}}

在 Python 中，安装第三方模块，是通过包管理工具 pip 完成的。

如果你正在使用 Mac 或 Linux，安装 pip 本身这个步骤就可以跳过了。

如果你正在使用
Windows，请参考\href{https://www.liaoxuefeng.com/wiki/1016959663602400/1016959856222624}{安装
Python}
一节的内容，确保安装时勾选了\texttt{pip}和\texttt{Add\ python.exe\ to\ Path}。

在命令提示符窗口下尝试运行\texttt{pip}，如果 Windows
提示未找到命令，可以重新运行安装程序添加\texttt{pip}。

注意：Mac 或 Linux 上有可能并存 Python 3.x 和 Python 2.x，因此对应的 pip
命令是\texttt{pip3}。

例如，我们要安装一个第三方库------Python Imaging Library，这是 Python
下非常强大的处理图像的工具库。不过，PIL 目前只支持到 Python
2.7，并且有年头没有更新了，因此，基于 PIL 的 Pillow
项目开发非常活跃，并且支持最新的 Python 3。

一般来说，第三方库都会在 Python 官方的
\href{https://pypi.python.org/}{pypi.python.org}
网站注册，要安装一个第三方库，必须先知道该库的名称，可以在官网或者 pypi
上搜索，比如 Pillow 的名称叫
\href{https://pypi.python.org/pypi/Pillow/}{Pillow}，因此，安装 Pillow
的命令就是：

\begin{pythoncode}
pip install Pillow
\end{pythoncode}

耐心等待下载并安装后，就可以使用 Pillow 了。

 
 \begin{figure}[htp]
	\centering
	\includegraphics[width=0.6\linewidth]{fig/1183096336546240l.png}
\end{figure}


\hypertarget{ux5b89ux88c5ux5e38ux7528ux6a21ux5757}{%
\subsubsection{安装常用模块}\label{ux5b89ux88c5ux5e38ux7528ux6a21ux5757}}

在使用 Python 时，我们经常需要用到很多第三方库，例如，上面提到的
Pillow，以及 MySQL 驱动程序，Web 框架 Flask，科学计算 Numpy 等。用 pip
一个一个安装费时费力，还需要考虑兼容性。我们推荐直接使用
\href{https://www.anaconda.com/}{Anaconda}，这是一个基于 Python
的数据处理和科学计算平台，它已经内置了许多非常有用的第三方库，我们装上
Anaconda，就相当于把数十个第三方模块自动安装好了，非常简单易用。

可以从 \href{https://www.anaconda.com/download/}{Anaconda 官网}下载 GUI
安装包，安装包有
500\textasciitilde600M，所以需要耐心等待下载。下载后直接安装，Anaconda
会把系统 Path 中的 python 指向自己自带的 Python，并且，Anaconda
安装的第三方模块会安装在 Anaconda 自己的路径下，不影响系统已安装的
Python 目录。

安装好 Anaconda 后，重新打开命令行窗口，输入 python，可以看到 Anaconda
的信息：

\begin{pythoncode}
┌────────────────────────────────────────────────────────┐
│Command Prompt - python                           - □ x │
├────────────────────────────────────────────────────────┤
│Microsoft Windows [Version 10.0.0]                      │
│(c) 2015 Microsoft Corporation. All rights reserved.    │
│                                                        │
│C:\> python                                             │
│Python 3.6.3 |Anaconda, Inc.| ... on win32              │
│Type "help", ... for more information.                  │
│>>> import numpy                                        │
│>>> _                                                   │
│                                                        │
│                                                        │
│                                                        │
└────────────────────────────────────────────────────────┘
\end{pythoncode}

可以尝试直接\texttt{import\ numpy}等已安装的第三方模块。

\hypertarget{ux6a21ux5757ux641cux7d22ux8defux5f84}{%
\subsubsection{模块搜索路径}\label{ux6a21ux5757ux641cux7d22ux8defux5f84}}

当我们试图加载一个模块时，Python 会在指定的路径下搜索对应的. py
文件，如果找不到，就会报错：

\begin{pythoncode}
>>> import mymodule
Traceback (most recent call last):
  File "<stdin>", line 1, in <module>
ImportError: No module named mymodule
\end{pythoncode}

默认情况下，Python
解释器会搜索当前目录、所有已安装的内置模块和第三方模块，搜索路径存放在\texttt{sys}模块的\texttt{path}变量中：

\begin{pythoncode}
>>> import sys
>>> sys.path
['', '/Library/Frameworks/Python.framework/Versions/3.6/lib/python36.zip', '/Library/Frameworks/Python.framework/Versions/3.6/lib/python3.6', ..., '/Library/Frameworks/Python.framework/Versions/3.6/lib/python3.6/site-packages']
\end{pythoncode}

如果我们要添加自己的搜索目录，有两种方法：

一是直接修改\texttt{sys.path}，添加要搜索的目录：

\begin{pythoncode}
>>> import sys
>>> sys.path.append('/Users/michael/my_py_scripts')
\end{pythoncode}

这种方法是在运行时修改，运行结束后失效。

第二种方法是设置环境变量\texttt{PYTHONPATH}，该环境变量的内容会被自动添加到模块搜索路径中。设置方式与设置
Path 环境变量类似。注意只需要添加你自己的搜索路径，Python
自己本身的搜索路径不受影响。

